\documentclass[a4paper,12pt]{exam}
\usepackage[T1]{fontenc}
\usepackage{lmodern}
\usepackage{comment}
\usepackage{fancyvrb}
\usepackage{graphicx}
\usepackage{geometry}
\usepackage{amssymb,amsmath}
\usepackage{ifxetex,ifluatex}
\usepackage{fixltx2e} % provides \textsubscript
% use upquote if available, for straight quotes in verbatim environments
\IfFileExists{upquote.sty}{\usepackage{upquote}}{}
\ifnum 0\ifxetex 1\fi\ifluatex 1\fi=0 % if pdftex
  \usepackage[utf8]{inputenc}
\else % if luatex or xelatex
  \ifxetex
    \usepackage{mathspec}
    \usepackage{xltxtra,xunicode}
  \else
    \usepackage{fontspec}
  \fi
  \defaultfontfeatures{Mapping=tex-text,Scale=MatchLowercase}
  \newcommand{\euro}{€}
\fi
% use microtype if available
\IfFileExists{microtype.sty}{\usepackage{microtype}}{}
\usepackage{longtable,booktabs}
\usepackage{graphicx}
% Redefine \includegraphics so that, unless explicit options are
% given, the image width will not exceed the width of the page.
% Images get their normal width if they fit onto the page, but
% are scaled down if they would overflow the margins.
\makeatletter
\def\ScaleIfNeeded{%
  \ifdim\Gin@nat@width>\linewidth
    \linewidth
  \else
    \Gin@nat@width
  \fi
}
\makeatother
\let\Oldincludegraphics\includegraphics
{%
 \catcode`\@=11\relax%
 \gdef\includegraphics{\@ifnextchar[{\Oldincludegraphics}{\Oldincludegraphics[width=\ScaleIfNeeded]}}%
}%
\ifxetex
  \usepackage[setpagesize=false, % page size defined by xetex
              unicode=false, % unicode breaks when used with xetex
              xetex]{hyperref}
\else
  \usepackage[unicode=true]{hyperref}
\fi
\hypersetup{breaklinks=true,
            bookmarks=true,
            pdfauthor={},
            pdftitle={Worksheet: Streams},
            colorlinks=true,
            citecolor=blue,
            urlcolor=blue,
            linkcolor=magenta,
            pdfborder={0 0 0}}
\urlstyle{same}  % don't use monospace font for urls
\setlength{\parindent}{0pt}
\setlength{\parskip}{6pt plus 2pt minus 1pt}
\setlength{\emergencystretch}{3em}  % prevent overfull lines
%\setcounter{secnumdepth}{0}

\geometry{
  %includeheadfoot,
  margin=2.54cm
}

\title{Exercises --- Week Seven}
\date{Spring term 2016}
\author{Optional exercise --- Object-Oriented Sets}

\begin{document}
\maketitle

You need to have installed the current Scala distribution and the current Java distribution
before commencing these exercises.
Don't forget the documentation for both distributions as they will come in very useful.

Any source code examples are available under the \verb!week07/objsets! folder.

The exercise is based upon one from one of the EPFL courses on functional programming techniques, 
later reused in Martin Odersky's \url{coursera.org} course.

\section*{Preamble}

In this exercise set you will work with an object-oriented representations based on binary trees.
You will be  completing the \verb!TweetSet.scala! file. 
This file defines an abstract class \verb!TweetSet! with two concrete subclasses, 
\begin{description}
\item[\texttt{Empty}] --- which represents an empty set, and 
\item[\texttt{NonEmpty(elem: Tweet, left: TweetSet, right: TweetSet)}] 
--- which represents a non-empty set as a binary tree rooted at \verb!elem!. 
\end{description}
The tweets are indexed by their text bodies: the bodies of all tweets on the left 
are lexicographically smaller than \verb!elem! and all bodies of elements on the right 
are lexicographically greater. Note that these classes are immutable: 
the set-theoretic operations do not modify \verb!this! but should return a new set.

Before tackling this assignment, we suggest you first study the already implemented 
methods of \verb!TweetSet!, \verb!contains! and \verb!incl!.
\begin{questions}
\question
Implement \emph{filtering} on tweet sets. Complete the stubs for the methods 
\verb!filter! and \verb!filterAcc!. \verb!filter! takes as argument a function, 
the predicate, which takes a tweet and returns a boolean. \verb!filter! 
then returns the subset of all the tweets in the original set for which the predicate is true. 
For example, the following call:
\begin{Verbatim}
tweets.filter(tweet => tweet.retweets > 10)	
\end{Verbatim}
applied to a set \verb!tweets! of two tweets, say, where the first tweet was not retweeted and 
the second tweet was retweeted 20 times should return a set containing only the second tweet.

Hint: start by defining the helper method \verb!filterAcc! 
which takes an accumulator set as a second argument. 
This accumulator contains the ongoing result of the filtering.
\begin{Verbatim}
/** This method takes a predicate and returns a subset of all the elements
 *  in the original set for which the predicate is true.
 */
def filter(p: Tweet => Boolean): TweetSet
def filterAcc(p: Tweet => Boolean, acc: TweetSet): TweetSet	
\end{Verbatim}
The definition of \verb!filter! in terms of \verb!filterAcc! should then be straightforward.

\question
Implement union on tweet sets. Complete the stub for the method \verb!union!. 
The method \verb!union! takes another set \verb!that!, and computes a new set which is 
the union of \verb!this! and \verb!that!, i.e., a set that contains exactly the elements 
that are either in \verb!this! or in \verb!that!, or \emph{in both}.
\begin{Verbatim}
def union(that: TweetSet): TweetSet	
\end{Verbatim}
Note that in this exercise it is your task to find out in which class(es) to define the 
\verb!union! method (should it be abstract in class \verb!TweetSet!?).

\question
The more often a tweet is ``re-tweeted'' 
(that is, repeated by a different user with or without additions), the more influential it is.

The goal of this part of the exercise is to add a method 
\verb!descendingByRetweet! to \verb!TweetSet! which should produce a linear sequence of 
tweets (as an instance of class \verb!TweetList!), ordered by their number of retweets:
\begin{Verbatim}
def descendingByRetweet: TweetList	
\end{Verbatim}
This method reflects a common pattern when transforming data structures. 
While traversing one data structure (in this case, a \verb!TweetSet!), 
we’re building a second data structure (here, an instance of class \verb!TweetList!). 
The idea is to start with the empty list \verb!Nil! (containing no tweets), 
and to find the tweet with the most retweets in the input \verb!TweetSet!. 
This tweet is removed from the \verb!TweetSet! (that is, we obtain a new \verb!TweetSet! 
that has all the tweets of the original set except for the tweet that was ``removed''; 
this \emph{immutable} set operation, \verb!remove!, 
is already implemented for you), and added to the result list by creating a new \verb!Cons!. 
After that, the process repeats itself, but now we are searching through a \verb!TweetSet! 
with one less tweet.

Hint: start by implementing the method \verb!mostRetweeted! which returns the most 
popular tweet of a \verb!TweetSet!.

\question
In the last step of this exercise your task is to detect influential tweets in a set of 
recent tweets. You are provided you with a \verb!TweetSet! containing several hundred tweets 
from popular tech news sites, located in the \verb!TweetReader! object (file 
\verb!TweetReader.scala!). \verb!TweetReader.allTweets! returns an instance of 
\verb!TweetSet! containing a set of all available tweets.

Furthermore, you are given two lists of keywords. The first list corresponds to 
keywords associated with Google and Android smartphones, while the second list corresponds 
to keywords associated with Apple and iOS devices. Your objective is to detect which 
platform has generated more interest or activity.

As a first step, use the functionality you implemented in the first parts of this assignment 
to create two different \verb!TweetSet!s, \verb!googleTweets! and \verb!appleTweets!. 
The first \verb!TweetSet!, \verb!googleTweets!, should contain all tweets that 
mention (in their ``text'') one of the keywords in the google list. 
The second \verb!TweetSet!, \verb!appleTweets!, should contain all tweets that mention 
one of the keyword in the apple list. Their signature is as follows:
\begin{Verbatim}
val googleTweets: TweetSet
val appleTweets: TweetSet	
\end{Verbatim}
Hint: use the \verb!exists! method of \verb!List! and \verb!contains! 
method of class \verb!java.lang.String!.

From the \emph{union} of those two \verb!TweetSet!s, produce \verb!trending!, 
an instance of class \verb!TweetList! representing a sequence of tweets ordered 
by their number of retweets:
\begin{Verbatim}
val trending: TweetList	
\end{Verbatim}

\end{questions}
\end{document}